\begin{longtable}{|p{3.5cm}|p{5.5cm}|p{6cm}|}

\hline

Categoria & Elementos, diretrizes e princípios identificados na literatura & Exemplos de estudos \\ \hline

\endfirsthead

\hline

Categoria & Elementos, diretrizes e princípios identificados na literatura & Exemplos de estudos \\ \hline

\endhead

\hline

\endfoot

\endlastfoot

Simplicidade e clareza da interface &

Linguagem simples e clara, textos concisos, redução de informações irrelevantes, redução da quantidade de opções apresentadas, organização clara das informações, interfaces simples e redução da complexidade visual. &

Breaking Barriers: The Design and Development of an Assistive Technology Web App for Older Latinos with Disabilities in Daily Activities; Inclusive websites for the elderly: User friendly guidelines for designers and managers of websites and applications; A study of web usability for older adults seeking online health resources. \\ \hline

Tamanho, legibilidade e apresentação visual &

Fontes grandes, texto legível, alto contraste, espaçamento entre palavras, botões grandes, ícones maiores e distintos, requisitos de apresentação visual e organização adequada do espaço em branco. &

Breaking Barriers: The Design and Development of an Assistive Technology Web App for Older Latinos with Disabilities in Daily Activities; Old and Young Users' White Space Preferences for Online News Web Pages; Usability Testing by Older Americans of a Prototype Google Map Web Site To Select Nursing Homes. \\ \hline

Navegação e orientação &

Navegação clara, hierarquia de informação, breadcrumbs, link para a página inicial, sitemap, mecanismo de busca, indicação de rolagem, retorno fácil à página inicial, sequência de ações e redução da profundidade da navegação. &

A study on the acceptance of website interaction aids by older adults; A usability study of websites for older travelers; Inclusive websites for the elderly: User friendly guidelines for designers and managers of websites and applications. \\ \hline

Consistência e previsibilidade &

Posicionamento consistente dos elementos, layout consistente, uso consistente de botões, terminologia adequada, estrutura previsível e manutenção de padrões de interação entre as páginas. &

Usability Testing by Older Americans of a Prototype Google Map Web Site To Select Nursing Homes; Inclusive websites for the elderly: User friendly guidelines for designers and managers of websites and applications; A methodological approach for designing and developing web-based inventories of mobile Assistive Technology applications. \\ \hline

Controles, botões e elementos interativos &

Botões grandes, ícones simples e claros, ícones de home, setas de voltar, controles de navegação, opções limitadas por página, áreas de interação maiores e posicionamento adequado dos botões. &

Breaking Barriers: The Design and Development of an Assistive Technology Web App for Older Latinos with Disabilities in Daily Activities; Usability Study to Promote Co-Creation among People with Disabilities, Developers, and Makers with a Focus on the Assistive Technology Open Platform in Korea; Designing Inclusive Smartwatch Interfaces: Guidelines for Enhancing Usability and Adoption Among Older Adults. \\ \hline

Feedback e prevenção de erros &

Feedback claro durante o processamento, mensagens de erro simples, prevenção de erros, mecanismos de recuperação, explicações para campos de formulários e indicação clara das ações realizadas pelo sistema. &

Usability Testing by Older Americans of a Prototype Google Map Web Site To Select Nursing Homes; A study on the acceptance of website interaction aids by older adults; Recoverability Walkthrough: An Alternative to Evaluate Digital Inclusion Interfaces. \\ \hline

Formulários e entrada de dados &

Formulários simplificados, quebra de formulários extensos em etapas, explicação dos campos, campos organizados, orientação durante o preenchimento e interfaces adaptadas às necessidades dos usuários. &

A Framework for Serving Inclusive Web Forms to Disabled and Elderly Individuals; A study on the acceptance of website interaction aids by older adults; Family-centered design: Interactive performance testing and user interface evaluation of the Slovenian eDavki public tax portal. \\ \hline

Multimodalidade e recursos alternativos &

Feedback multissensorial, entrada multimodal, suporte por voz, explicação falada, informações alternativas para áudio e vídeo, imagens, fotos e vídeos como recursos complementares. &

Designing Inclusive Smartwatch Interfaces: Guidelines for Enhancing Usability and Adoption Among Older Adults; U-Cker: Initial Development of an Interactive Video Retrieval System for Novice Users; Content Analysis on User-Focused Support Features of Health Support Sites for Geriatric Depression. \\ \hline

Personalização e adaptação da interface &

Personalização da interface, adaptação às características dos usuários, preferências individuais, recursos ajustáveis e modelos de usuário para adaptar a interação às necessidades específicas. &

Interface personalization through inclusive user modelling; Conducting acceptance tests for elderly people on the web: Using the gpii preference set for a personalized evaluation; Patterns for user interface adaptations towards runtime adaptations for improving the usability of web forms for elderly users. \\ \hline

Acessibilidade e inclusão &

Conformidade com WCAG, alto contraste, suporte a tecnologias assistivas, navegação por teclado, acessibilidade para diferentes deficiências, reorganização do conteúdo e mecanismos alternativos de entrada e interação. &

A methodological approach for designing and developing web-based inventories of mobile Assistive Technology applications; Breaking Barriers: The Design and Development of an Assistive Technology Web App for Older Latinos with Disabilities in Daily Activities; Research on Web Interface barrier-free Design for Elderly. \\ \hline

\caption{Elementos de interface, diretrizes e princípios de design identificados nos estudos que respondem diretamente à RQ2.}

\label{tab:elementosRQ2}

\end{longtable}